\section{Discussion}
\label{sec:discussion}

The results come from two kinds of test, one causal and one observational. H2's causal result shows contamination as a training-time memorisation effect: the fine-tuned checkpoint reproduces its training targets verbatim and passes near-perfectly on the exact problems it trained on. H1's observational results are weaker, and they ask a different question. Do models trained by others, on corpora nobody here can inspect, already carry this effect for ULT's published content? The two results do not conflict. A causal probe with one controlled intervention can show a large, clear effect. An observational comparison across nine models with unknown training histories cannot be expected to recover the same effect as clearly.

One further finding is worth stating. After fine-tuning, the mutation score on the seen-ULT split is $\HTwoPostMutk{}$, not \MutkCeiling{}, even though the model reproduces its training target word for word. Byte-identical recall of a reference test suite does not guarantee the recalled suite kills every mutant the scoring harness generates. Memorisation alone can only inflate a mutation-testing score so far, whether or not the benchmark resists contamination.

The benchmark's validity argument is top-down by design: complexity tiers, mutation operators, and scoring metric are all specified in advance from an explicit model of what test-generation capability should look like. H6 supplies the complementary bottom-up check: Does the benchmark's empirical ranking of models agree with an independent capability measure? The answer is mixed. The conditional primary metric lands inside the pre-registered validity band, though it does not reach significance at this sample size. The coverage-adjusted co-primary exceeds the band. Breaking that figure down shows most of its agreement is about executability, not about mutation-testing quality. Neither reading settles the question alone. Together they place DyMuT close to its intended target, correlated with general code-generation ability.

Several threats to validity apply. Running untrusted code in a subprocess is not the same as running it safely. An early ground-truth quality check executed thousands of real-world ULT functions with the working directory left at the project root. Several of those functions wrote files as a side effect of ordinary relative-path I/O. The check now runs from an isolated temporary directory. The causal fine-tuning result (H2) uses a single model at a single, small parameter scale; whether the same clean separation between memorisation and transfer holds at larger scale is untested here. Mutation score is conditional on baseline correctness throughout. Section~\ref{sec:results} shows this affects cross-benchmark rank agreement in H6. It is a property of the metric definition, not of the benchmark itself. ULT's own published ground truth passes this study's harness at a lower rate than its role as a correctness reference would suggest, a finding supported by an internal recomputation of the benchmark's own complexity labels and by a direct contradiction between the source paper's stated mutation-score formula and its own published scoring code: this study's co-primary (coverage-adjusted) definition matches the source paper's prose formula, while this study's primary (conditional) definition matches the paper's own published code instead, which is why both are reported as co-primaries throughout. None of this study's reported metrics are affected by the contradiction, since it is disclosed and accounted for; the finding is reported here as a caveat on ULT's suitability as a ground-truth reference.
